\section{Formal verification in Lean 4}
\label{sec:lean}

The results listed below have been formalized in Lean~4 with Mathlib.
The integrated release contains 28 source modules and passed a clean project
\texttt{lake build} with Lean~4.36.0-rc2 and Mathlib revision
\texttt{9d581dfff2e41de989072ea63abf3daf69311520}.
The dependency source trees were checked against the pinned manifest;
existing dependency build caches were reused. All 112 named declarations in
the integrated audit passed the stated axiom check. Sources, build records,
and SHA-256 fingerprints are provided in the companion repository
\url{https://github.com/CWBaker360/nseries-lean-verification}.
The repository commit identifies the complete release snapshot.
Table~\ref{tab:lean} lists the correspondence between the statements of this
note and the Lean theorems.

\begin{table}[h]
\centering
\begin{tabular}{@{}p{0.43\linewidth}p{0.51\linewidth}@{}}
\toprule
Statement & Lean theorem \\
\midrule
Lemma 2.2 (even-power expansion) & \texttt{nseries\_sine\_remainder\_bound} \\
Theorem 3.1 (two-scale, general sequence) & \texttt{two\_scale\_general\_bound} \\
Corollary 3.2 (polygonal two-scale) & \texttt{richardson\_error\_bound} \\
Theorem 4.2 ($\pi_g^{(k)}=\pi$) & \texttt{acceleratedPi\_eq\_pi} \\
Theorem 6.1 (polygonal specialization) & \texttt{three\_scale\_error\_bound}, \texttt{three\_scale\_sharp} \\
Corollary 6.2 (analytic dyadic form) & \texttt{dyadic\_three\_scale\_sharp} \\
Theorem 7.1 (moment cancellation) & \texttt{general\_moment\_cancellation} \\
Remark 7.2 (uniqueness of weights) & \texttt{moment\_weights\_unique} \\
Moment-weight existence & \texttt{nseries\_weights\_existsUnique} \\
Explicit Lagrange weights & \texttt{lagrange\_weights\_moments}, \texttt{lagrange\_weight\_formula} \\
Section 10 (trisection cubic) & \texttt{chordHalf\_trisection\_unique} \\
Polygonal limit & \texttt{polygonalPi\_eq\_pi} \\
Corollary 3.2 on intrinsic polygons & \texttt{intrinsic\_two\_scale\_bound} \\
Corollary 6.2 on intrinsic polygons & \texttt{intrinsic\_dyadic\_three\_scale\_sharp} \\
\bottomrule
\end{tabular}
\caption{Correspondence between this note and the Lean development.}
\label{tab:lean}
\end{table}

The two named three-scale bounds in the table apply to the polygonal analytic sequence; the generic finite cancellation mechanism is covered by the explicit-hypothesis moment theorem.

Four points of scope should be stated precisely.

\emph{Explicit bounds.} The asymptotic statements are verified in
explicit-constant form. For example, Lemma~2.2 is proved as
\[
\left|N\sin\frac{\pi}{N}
-\Bigl(\pi-\frac{\pi^3}{6N^2}+\frac{\pi^5}{120N^4}-\frac{\pi^7}{5040N^6}\Bigr)\right|
\le \frac{\pi^9}{362880\,N^8}
\qquad (N>0),
\]
and the three-scale operator is shown to satisfy
\[
\widehat P_N^{(k,3)}=\pi-\frac{\pi^7}{5040\,k^6N^6}+O(N^{-8})
\]
with an explicit constant. Theorems~3.1 and~7.1 are formalized under an
explicit remainder hypothesis $|q(M)-Q-\cdots|\le B/M^{p}$ rather than a
big-$O$ hypothesis; the Lean statements are therefore quantitative versions
of the theorems as stated here, not verbatim transcriptions.

\emph{The polygonal limit.} The half-perimeter of the regular $n$-gon is
defined in Lean without reference to $\pi$, as
\[
\frac{n}{2}\min_{\substack{z^n=1\\ z\ne 1}}|1-z| \qquad (n\ge 3),
\]
and the limit of this sequence is proved to equal $\pi$. The definition is
$\pi$-free; the proof of the identification uses the classical constant.
This is the raw polygonal limit, distinct from the accelerated limit
$\pi_g^{(k)}$ of Definition~4.1, which is verified separately.
The acceleration bounds are proved for the analytic sequence
$N\sin(\pi/N)$, and the intrinsic half-perimeter is proved equal to that
sequence for $n\ge 3$. Two of the bounds are also stated and verified
directly for the intrinsic sequence: the two-scale bound of Corollary~3.2
for integer $k\ge 2$, and the sharp dyadic three-scale bound of
Corollary~6.2. The remaining bounds (the general-$k$ three-scale bound and
the accelerated limit of Theorem~4.2) have not been restated for the
intrinsic sequence.

\emph{What is not verified.} The digit tables of Section~10, the
$100{,}000$-digit agreement of the two implementations, and the observed
growth of the number of correct digits with the number of levels are
computational evidence, not formally verified theorems. For the trisection
recurrence, Lean verifies the cubic identity and the uniqueness of its root
in $[0,\tfrac12]$ for $N\ge 2$, not the numerical runs.

\emph{Weights.} The Vandermonde formalization establishes existence and
uniqueness of the moment weights at distinct nodes. Separately, the explicit
Lagrange-basis weights
\[
w_i=\prod_{k\ne i}\frac{-x_k}{x_i-x_k}
\]
are proved to satisfy the moment system. These results certify each finite
moment-cancellation system, not a uniform digit-growth rate as the number of
levels increases.

The reported endpoint axiom audits list only \texttt{propext}, \texttt{Classical.choice},
and \texttt{Quot.sound}; no \texttt{sorry}, \texttt{admit}, or additional axioms were reported
in the audited modules.
